\newpage
\section{问题重述}

\subsection{问题背景}
智慧园区是城市能源转型的重要单元，其能耗管理直接关系到碳排放目标的实现。园区内通常包含建筑负荷、分布式光伏、储能系统和充电桩等多种能源设施，如何实现多能协同优化调度是智慧园区运营的核心问题。

\subsection{问题提出}
\textbf{问题一}：对园区历史能耗数据进行时空特征分析，识别各类负荷的用能规律和能耗强度，建立能耗分解模型。

\textbf{问题二}：建立园区微电网多能互补调度模型，以运行成本最小为目标，考虑光伏出力不确定性，设计鲁棒调度策略。

\textbf{问题三}：引入碳交易机制，建立考虑碳排放成本的多目标调度模型，分析不同碳价情景下的调度策略变化。

\subsection{研究意义}
为园区级综合能源系统优化调度提供理论方法，助力碳达峰碳中和目标实现。
